\begin{textsample}{POS Dim 1 – vem\_ed\_06 – Score 12.00 – vem\_ed\_06/t024.txt}  \label{ex:f1_pos_128}
Presença e energia Neste primeiro semestre de 2021, ocorre a \textbf{quinta} edição do Projeto Colaborativo Internacional (PCI) entre Fatecs e Tianjin Normal University (China), focado em \textbf{competências} interculturais e \textbf{linguísticas} (aprendizagem de inglês por falantes não nativos).
Participam as unidades de Franca, Campinas, Catanduva, São José do Rio Preto, Itaquaquecetuba, Piracicaba, Sebrae (na capital) e Mogi das Cruzes.
Neusa Haruka Sezaki Gritti, envolvida desde o \textbf{início} do projeto, relata sua evolução.
A \textbf{cada} edição a qualidade do \textbf{processo} e os resultados melhoram.
No segundo semestre de 2020, o \textbf{número} de mensagens entre os estudantes brasileiros e os chineses no semestre \textbf{foi} o maior registrado nas cinco edições desse Intercâmbio Virtual: 3080.
Uma das fases \textbf{que} os alunos mais curtem \textbf{é} o icebreaker (quebra-gelo).
Os participantes compreendem \textbf{que}, além da comunicação em língua estrangeira, há outras barreiras: a cultura de \textbf{cada} país, a pesquisa para explorar os assuntos propostos, além da dificuldade de se entender o accent (sotaque) do estrangeiro numa reunião online \textbf{síncrona}.
Essa atividade traz um pouco de insegurança para os alunos brasileiros, mas num instante eles se engajam e se divertem muito.
Além do envolvimento cultural, \textbf{linguístico} e tecnológico - objetivos do PCI- a satisfação nesse trabalho \textbf{é} receber feedbacks de alunos comentando \textbf{que}"\textbf{foi} uma das melhores experiências do semestre, indicaria para \textbf{todos}";"as \textbf{tarefas} não me ajudaram somente a melhorar meu inglês, mas a mim mesmo";"gostei de trabalhar com os chineses, e decidi \textbf{aprender} chinês".
Mas, para \textbf{que} um Projeto Colaborativo Internacional \textbf{seja} \textbf{bem-sucedido}, há vários \textbf{pontos} \textbf{que} precisam \textbf{ser} \textbf{destacados}: 1.Parceria.
Os dois professores (brasileiro e estrangeiro) precisam estar em sintonia sempre e deve haver muito respeito de ambos, não interferindo na atitude do outro junto aos estudantes; 2.Selecionar um ótimo líder para \textbf{cada} grupo; 3.
Pelo menos um membro do grupo deve ter fluência na língua estrangeira; 4.
O professor precisa ficar bastante atento a tudo \textbf{que} acontece, digo, nas mensagens de \textbf{todos} os grupos e interferir sempre \textbf{que} necessário; 5.
Ficar atento também às soft skills necessárias para um bom relacionamento e excelentes resultados: comunicação; trabalho em equipe, resiliência, administração de tempo, liderança, \textbf{disciplina} e organização.
E o mais importante: o professor precisa ter muita energia, ânimo e estar sempre presente com \textbf{todos}.
Quando há um problema com algum aluno, auxiliá-lo para \textbf{que} esteja sempre engajado no \textbf{processo}.

% matched lemmas: aprender, bem-sucedido, cada, competência, destacar, disciplina, início, linguístico, número, ponto, processo, que, quinto, ser, síncrono, tarefa, todo
\end{textsample}
